The conditional comparisons in the preliminary mixing argument use
Ginibre's formulation of ferromagnetic correlation inequalities and
the positive-association theorem of Fortuin, Kasteleyn and Ginibre~\cite{Ginibre1970,FortuinKasteleynGinibre1971}.
The passage from local disagreement estimates to loss of boundary
influence is related to disagreement percolation~\cite{VanDenBergMaes1994}.
In exact blocking, the connected polymer sums use the classical
tree-graph approach to cluster expansions; a useful formulation of
the Penrose identity is given by Fern\'andez and
Procacci~\cite{FernandezProcacci}.

The finite-volume comparison uses reflection positivity and the
chessboard method systematized by Fr\"ohlich, Israel, Lieb and
Simon~\cite{FILS78,FILS80}.  For the final passage from Euclidean
distributions to a local relativistic field, we use the reconstruction
framework of Osterwalder and Schrader~\cite{OS73} in the corrected
form with the linear growth condition of \cite[Section~IV.1]{OS75}.
Reflection positivity, Euclidean covariance, clustering and the
hierarchy of distributional bounds are consequently separate parts
of the continuum argument.
